\section{Substituição de fórmulas}\label{sec:substituicaoDeFormulas}

Intuitivamente, se uma teoria prova que duas fórmulas são equivalentes, então elas podem ser substituídas uma pela outra em qualquer contexto.

Nessa seção, formalizaremos esse raciocínio.

Lembremos que, na Seção~\ref{sec:teoriasPrimeiraOrdem}, definimos a extensão da lógica proposicional para a lógica de primeira ordem $T_P$, e a substituição de letras proposicionais por fórmulas de uma linguagem de primeira ordem.

\begin{metatheorem}\label{mthm:substituicaoFormulasEquivalentes}
    Seja $T$ uma teoria de primeira ordem.
    Sejam $\varphi_1,\dots,\varphi_n$ e $\psi_1,\dots,\psi_n$ fórmulas da linguagem de $T$, com $n\geq 1$, tais que, para cada $i\in\{1,\dots,n\}$
    \[
        T\vdash \varphi_i\leftrightarrow\psi_i.
    \]

    Considere, na linguagem $\mathcal L_P$, letras proposicionais distintas $A_1,\dots,A_n$ e uma $\mathcal L_P$-fórmula $\theta$ em que todas as letras proposicionais que ocorrem em $\theta$ estão entre $A_1,\dots,A_n$.

    Sejam $\rho_\varphi$ e $\rho_\psi$ as atribuições dadas por
    \[
        \rho_\varphi(A_i)=\varphi_i
        \qquad\text{e}\qquad
        \rho_\psi(A_i)=\psi_i
    \]
    para cada $i\in\{1,\dots,n\}$.

    Então
    \[
        T\vdash \subst(\theta,\rho_\varphi)\leftrightarrow \subst(\theta,\rho_\psi).
    \]
\end{metatheorem}

\begin{proof}
    Procedemos por indução na complexidade de $\theta$.
    Para cada subfórmula $\eta$ de $\theta$, escreva
    \[
        \eta^\varphi=\subst(\eta,\rho_\varphi)
        \qquad\text{e}\qquad
        \eta^\psi=\subst(\eta,\rho_\psi).
    \]
    Mostraremos que, para toda subfórmula $\eta$ de $\theta$,
    \[
        T\vdash \eta^\varphi\leftrightarrow \eta^\psi.
    \]

    Se $\eta$ é uma letra proposicional $A_i$, então $\eta^\varphi=\varphi_i$ e $\eta^\psi=\psi_i$.
    Como $T\vdash \varphi_i\leftrightarrow\psi_i$, segue a tese.

    Se $\eta$ é uma fórmula equacional $x=y$, então $\eta^\varphi=\eta^\psi=x=y$.
    Portanto, $T\vdash \eta^\varphi\leftrightarrow\eta^\psi$, pois $(x=y)\leftrightarrow(x=y)$ é uma instância de tautologia.

    Se $\eta$ é $\neg\chi$, então, por hipótese de indução, $T\vdash \chi^\varphi\leftrightarrow\chi^\psi$.
    Pela preservação da negação por equivalências,
    \[
        T\vdash \neg\chi^\varphi\leftrightarrow\neg\chi^\psi.
    \]

    Se $\eta$ é $\chi\square\xi$, em que $\square$ é um conectivo binário, então, por hipótese de indução,
    \[
        T\vdash \chi^\varphi\leftrightarrow\chi^\psi
        \qquad\text{e}\qquad
        T\vdash \xi^\varphi\leftrightarrow\xi^\psi.
    \]
    Por raciocínio proposicional, isto é, pela instância de tautologia correspondente à preservação de $\square$ por equivalências,
    \[
        T\vdash
        (\chi^\varphi\square\xi^\varphi)
        \leftrightarrow
        (\chi^\psi\square\xi^\psi).
    \]

    Finalmente, suponha que $\eta$ é $Qx\chi$, em que $Q$ é $\forall$ ou $\exists$.
    Por hipótese de indução,
    \[
        T\vdash \chi^\varphi\leftrightarrow\chi^\psi.
    \]
    Se $Q$ é $\forall$, a preservação do quantificador universal por equivalências implica $T\vdash \forall x\chi^\varphi\leftrightarrow\forall x\chi^\psi$.

    Se $Q$ é $\exists$, a preservação do quantificador existencial por equivalências implica
    \[
        T\vdash \exists x\chi^\varphi\leftrightarrow\exists x\chi^\psi.
    \]
    Em ambos os casos,
    \[
        T\vdash \eta^\varphi\leftrightarrow\eta^\psi.
    \]

    Aplicando a conclusão à própria fórmula $\theta$, obtemos
    \[
        T\vdash \subst(\theta,\rho_\varphi)\leftrightarrow \subst(\theta,\rho_\psi).
    \]
\end{proof}

Nos exemplos abaixo, usaremos livremente que, para toda fórmula $\alpha$, temos $T\vdash \alpha\leftrightarrow\alpha$ por uma instância de tautologia.
Assim, para manter uma parte do contexto inalterada, basta atribuir a mesma fórmula à letra proposicional correspondente nas duas atribuições.

\begin{example}
    Suponha que $T\vdash \varphi\leftrightarrow\psi$.
    Tomando $\theta$ como $\neg A_1$, obtemos
    \[
        T\vdash \neg\varphi\leftrightarrow\neg\psi.
    \]
    Assim, fórmulas equivalentes podem ser substituídas dentro de negações.
\end{example}

\begin{example}
    Suponha que
    \[
        T\vdash \varphi\leftrightarrow\psi
        \qquad\text{e}\qquad
        T\vdash \chi\leftrightarrow\eta.
    \]
    Tomando $\theta$ como $(A_1\wedge A_2)\rightarrow A_1$, obtemos
    \[
        T\vdash
        \bigl((\varphi\wedge\chi)\rightarrow\varphi\bigr)
        \leftrightarrow
        \bigl((\psi\wedge\eta)\rightarrow\psi\bigr).
    \]
    Esse exemplo ilustra que várias substituições equivalentes podem ser feitas simultaneamente.
\end{example}

\begin{example}
    Suponha que $T\vdash \varphi\leftrightarrow\psi$.
    Tomando $\theta$ como $\forall x A_1$, obtemos
    \[
        T\vdash \forall x\varphi\leftrightarrow\forall x\psi.
    \]
    Analogamente, tomando $\theta$ como $\exists x A_1$, obtemos
    \[
        T\vdash \exists x\varphi\leftrightarrow\exists x\psi.
    \]
    Portanto, a substituição por fórmulas equivalentes também é preservada sob quantificadores.
\end{example}

\begin{example}
    Na linguagem da teoria dos conjuntos, suponha que
    \[
        T\vdash \varphi(z)\leftrightarrow\psi(z).
    \]
    Aplicando o Metateorema~\ref{mthm:substituicaoFormulasEquivalentes} ao esqueleto
    \[
        \theta \equiv \forall z(A_1\rightarrow A_2),
    \]
    com
    \[
        \rho_\varphi(A_1)=z\in a,\qquad
        \rho_\psi(A_1)=z\in a,
    \]
    e
    \[
        \rho_\varphi(A_2)=\varphi(z),\qquad
        \rho_\psi(A_2)=\psi(z),
    \]
    obtemos
    \[
        T\vdash
        \forall z(z\in a\rightarrow\varphi(z))
        \leftrightarrow
        \forall z(z\in a\rightarrow\psi(z)).
    \]
    Esse exemplo ilustra que podemos substituir apenas uma parte de uma fórmula por outra equivalente, deixando o restante do contexto inalterado ao atribuir a mesma fórmula à letra proposicional correspondente nos dois lados.
\end{example}

\begin{example}
    Suponha que, na linguagem da teoria dos conjuntos,
    \[
        T\vdash \varphi(z)\leftrightarrow\psi(z).
    \]
    Podemos substituir apenas uma das ocorrências de $\varphi(z)$ em uma fórmula em que ela aparece mais de uma vez.

    Por exemplo, considere a fórmula
    \[
        \forall z\bigl(z\in a\rightarrow(\varphi(z)\rightarrow\varphi(z)\wedge z\in b)\bigr).
    \]
    Queremos substituir apenas a segunda ocorrência de $\varphi(z)$ por $\psi(z)$, mantendo a primeira ocorrência inalterada.

    Para isso, tomemos o esqueleto
    \[
        \theta \equiv \forall z\bigl(A_1\rightarrow(A_2\rightarrow A_3\wedge A_4)\bigr),
    \]
    com
    \[
        \rho_\varphi(A_1)=z\in a,\qquad
        \rho_\varphi(A_2)=\varphi(z),\qquad
        \rho_\varphi(A_3)=\varphi(z),\qquad
        \rho_\varphi(A_4)=z\in b,
    \]
    e
    \[
        \rho_\psi(A_1)=z\in a,\qquad
        \rho_\psi(A_2)=\varphi(z),\qquad
        \rho_\psi(A_3)=\psi(z),\qquad
        \rho_\psi(A_4)=z\in b.
    \]
    Pelo Metateorema~\ref{mthm:substituicaoFormulasEquivalentes}, obtemos
    \[
        T\vdash
        \forall z\bigl(z\in a\rightarrow(\varphi(z)\rightarrow\varphi(z)\wedge z\in b)\bigr)
        \leftrightarrow
        \forall z\bigl(z\in a\rightarrow(\varphi(z)\rightarrow\psi(z)\wedge z\in b)\bigr).
    \]
    Assim, uma mesma fórmula pode ocorrer várias vezes em um contexto, e podemos substituir apenas algumas de suas ocorrências, tratando-as por letras proposicionais distintas no esqueleto.
\end{example}

\subsection*{Exercícios}

\begin{exercise}
    Suponha que $T\vdash \varphi\leftrightarrow\psi$.
    Use o Metateorema~\ref{mthm:substituicaoFormulasEquivalentes} para mostrar que:
    \begin{enumerate}
        \item $T\vdash \neg\varphi\leftrightarrow\neg\psi$;
        \item $T\vdash (\varphi\wedge\chi)\leftrightarrow(\psi\wedge\chi)$;
        \item $T\vdash (\chi\rightarrow\varphi)\leftrightarrow(\chi\rightarrow\psi)$;
        \item $T\vdash (\varphi\rightarrow\chi)\leftrightarrow(\psi\rightarrow\chi)$.
    \end{enumerate}
    Em cada item, explicite o esqueleto proposicional $\theta$ e as atribuições utilizadas.
\end{exercise}

\begin{exercise}
    Suponha que
    \[
        T\vdash \varphi\leftrightarrow\psi
        \qquad\text{e}\qquad
        T\vdash \chi\leftrightarrow\eta.
    \]
    Use o Metateorema~\ref{mthm:substituicaoFormulasEquivalentes} para concluir que
    \[
        T\vdash
        ((\varphi\vee\chi)\rightarrow(\varphi\wedge\chi))
        \leftrightarrow
        ((\psi\vee\eta)\rightarrow(\psi\wedge\eta)).
    \]
\end{exercise}

\begin{exercise}
    Suponha que $T\vdash \varphi\leftrightarrow\psi$.
    Conclua que
    \[
        T\vdash
        \neg(\chi\wedge(\varphi\vee\eta))
        \leftrightarrow
        \neg(\chi\wedge(\psi\vee\eta)).
    \]
\end{exercise}

\begin{exercise}
    Seja $T$ uma teoria de primeira ordem e suponha que
    \[
        T\vdash \varphi_i\leftrightarrow\psi_i
    \]
    para cada $i\in\{1,\dots,n\}$, com $n\geq 1$.
    Mostre, usando o Metateorema~\ref{mthm:substituicaoFormulasEquivalentes}, que
    \[
        T\vdash
        \bigwedge_{i=1}^n\varphi_i
        \leftrightarrow
        \bigwedge_{i=1}^n\psi_i.
    \]
\end{exercise}
